\section{边界条件与反例检验}

前文展示了重型 population training 的潜力，本节反向检查它何时不值得。判断依据不是团队是否喜欢多 Agent，而是任务开放度、失败代价、用户异质性和漏洞增长速度；只有这些因素超过单策略训练与静态回归的承载能力，league 的额外复杂度才有回报。

到这一步，一个自然问题是：既然 Multi-Agent 这么重要，是否所有 LLM 系统都应该立刻采用 league + exploiter 的重型方案？答案是否定的。工程上必须先识别边界条件，再决定投入强度。

\begin{figure}[H]
\centering
\includegraphics[width=0.92\textwidth]{frame-03.jpg}
\caption{讲者对 ``current LLMs'' 的局限性对照}
\label{fig:current-llm-limits}
\end{figure}

\begin{knowledgebox}{读图：当前 LLM 与 AlphaStar 的三处资源差异}
人类数据已经大量进入 LLM，reward 却只在少数可验证域完整；总体计算仍重 pretraining，post-training 相对轻；用户多为协作者但攻击持续存在。该图支持“后训练仍有扩展空间”，不证明把算力简单搬到 RL 就一定得到同等收益。
\end{knowledgebox}
图 \ref{fig:current-llm-limits} 展示了讲者在后半段对现状的判断\footnote{来源：\href{https://www.youtube.com/watch?v=ntjOxjZMaac}{Lecture 03 视频帧}，时间戳 00:30:55。}：奖励难以完备定义、计算分配偏向 pretraining、任务和数据覆盖仍有限。换句话说，问题并不是``要不要多 Agent''，而是``何时、以何种成本和治理方式引入''。

\subsection{哪些场景不应立即上重型 Multi-Agent}
如果任务满足以下特征，先做单 Agent 基础能力打磨更划算：任务边界固定、工具集合很小、错误成本可控、用户行为高度同质。这些条件下，引入多策略生态的收益可能不足以覆盖复杂度成本。

\begin{warningbox}{反例 1：小而确定的自动化任务}
例如固定格式报表生成、单 API 数据同步、只读检索问答。此时首要瓶颈通常是数据清洗、接口稳定性和监控，而不是对抗学习能力。过早引入 league 机制会显著增加维护负担。
\end{warningbox}

\begin{knowledgebox}{反例不等于否定}
即便暂不使用 multi-agent 训练，也应保留最小化的对抗评测接口。因为业务一旦扩展到开放用户交互，系统会迅速进入讲者描述的``协作+对抗共存''状态。
\end{knowledgebox}

\subsection{什么时候必须升级为 population 训练}
当系统出现以下信号时，应尽快引入 exploiter 与动态匹配：
\begin{enumerate}
    \item 线上 failure mode 的长尾持续增长，且人工规则修补速度跟不上。
    \item 模型在基准集表现稳定，但在真实用户多轮交互中回归明显。
    \item 同一漏洞在不同 prompt 变体下反复出现，表明是策略结构问题而非样本噪声。
\end{enumerate}

\begin{importantbox}{升级触发器}
升级的核心触发器不是``模型分数高低''，而是``失败发现速度是否超过修复速度''。一旦失败发现速度长期更快，说明系统已进入需要 population dynamics 的阶段。
\end{importantbox}

\subsection{治理与成本：课程之外的现实问题}
讲者谈的是研究范式，但产品侧还要面对治理和成本问题。Multi-Agent 系统往往引入更多执行链路、更多状态和更多权限边界，因此必须把安全治理与审计能力前置。

\begin{table}[H]
\centering
\caption{Multi-Agent 升级的成本收益审视（实践清单）}
\begin{tabular}{p{2.7cm}p{4.4cm}p{5.5cm}}
\toprule
维度 & 收益 & 成本/风险 \\
\midrule
鲁棒性 & 更快发现并修复长尾漏洞 & 训练与评测流水线复杂度上升 \\
泛化性 & 对新策略/新任务迁移更稳健 & 需要持续数据运营与任务调度 \\
安全性 & 可系统化 red teaming & 若权限设计不当，攻击面可能扩大 \\
可解释性 & 轨迹更完整，便于审计 & 追踪与存储成本显著增加 \\
\bottomrule
\end{tabular}
\end{table}

\begin{warningbox}{治理底线}
任何引入自动执行能力的代理系统，都必须具备最小回滚机制与高风险操作熔断机制。没有这两项，模型再强也不应进入高权限生产环境。
\end{warningbox}

\subsection{本章小结}
本章给出一个务实判断：Multi-Agent 不是默认答案，而是阶段性答案。先识别业务所处阶段，再决定训练复杂度；一旦进入开放交互与高风险执行场景，就应尽早升级到 population 视角，并同步建设治理能力。

